\chapter*{Abstract}
\addcontentsline{toc}{chapter}{Abstract}
\markboth{Abstract}{Abstract}

Laser-Speckled Background Oriented Schlieren (SBOS) is an optical density measurement technique that replaces the printed background of traditional BOS with a laser speckle pattern. This background has higher contrast, is tunable through the optics, and is always in focus, allowing unusual reference placements. This thesis develops a design framework and tests it in four experimental campaigns. The framework shows that negative defocus improves the sensitivity-resolution ratio over standard BOS. An angled glass plate validated its predictions. A laser speckle study produced speckle sizes of 2--5 px, optimal for cross-correlation. An under-expanded jet yielded a quantitative density field that matched literature data over the first shock cell. SBOS imaged a linear aerospike nozzle in a high-speed wind tunnel. With the focus inside the test section, it approached schlieren resolution, and a displacement-ratio method retrieved the sensitivity at zero defocus. This work lays the groundwork for SBOS in high-speed wind tunnels.
